\subsection{WebSocket 实时通信机制}

WebSocket 是一种基于 TCP 的全双工通信协议，由 Fette 和 Melnikov 于 2011 年标准化为 RFC 6455\footnote{Fette, I., Melnikov, A. “The WebSocket Protocol.” RFC 6455, IETF, 2011.}。与传统的 HTTP 通信采用“客户端请求、服务器响应”的模式不同，WebSocket 在客户端与服务器建立连接后，可以在同一连接上进行双向、全双工的实时数据交互，因此常用于在线聊天、实时通知、协同编辑和数据监控等场景。

WebSocket 通信的一般过程包括连接建立、身份认证、消息发送、消息转发和连接关闭五个阶段。在连接建立阶段，客户端首先向服务器发送 HTTP 升级请求（Upgrade: websocket），服务器返回 101 Switching Protocols 响应后双方建立 WebSocket 长连接。此后，客户端和服务器可以在该连接上持续交换数据帧（Data Frame），而不需要为每条消息重新建立 HTTP 请求。这一机制显著降低了通信延迟和带宽开销，适合需要频繁消息交互的实时场景。

在 FoodShield 系统中，WebSocket 用于实现用户端与骑手端之间的订单级实时通信。系统以订单号为基本单位建立通信房间，每个订单对应一个独立的通信通道。用户或骑手进入通信房间前，服务器需要验证其身份状态和订单参与资格。对于用户端，服务器需要校验订单号和 Token；对于骑手端，服务器需要校验其是否为该订单绑定的配送方。只有验证通过的参与方才能加入对应订单房间。

订单通信过程可以抽象为：

\begin{figure}[H]
\centering
\begin{tikzpicture}[>=Latex, font=\small]

% 顶部参与方
\node[font=\bfseries] (client) at (0,0) {客户端};
\node[font=\bfseries] (server) at (4,0) {服务器};
\node[font=\bfseries] (room)   at (8,0) {订单房间};

% 生命线
\draw[densely dashed] (0,-0.4) -- (0,-3.4);
\draw[densely dashed] (4,-0.4) -- (4,-3.4);
\draw[densely dashed] (8,-0.4) -- (8,-3.4);

% 交互过程
\draw[->] (0,-1.0) -- node[above] {发送 \textit{orderID, credential}} (4,-1.0);
\draw[->] (4,-1.9) -- node[above] {验证凭证并加入房间} (8,-1.9);
\draw[->] (4,-2.8) -- node[above] {转发消息} (8,-2.8);

% 房间内广播说明
\node[align=center] at (8,-3.7) {\small 广播至\\用户 / 骑手};

\end{tikzpicture}
\caption{订单认证与消息转发时序图}
\end{figure}

其中，$\mathit{credential}$表示订单认证凭证（即 Token）或骑手身份凭证。服务器在转发消息的同时，对消息进行加密存储、哈希计算和日志记录，使通信功能与安全审计机制形成联动。

% !TEX root = ../../main.tex
% WebSocket 握手时序图
\begin{figure}[H]
\centering
\begin{tikzpicture}[
  arr/.style   = {-Stealth, thick},
  lbl/.style   = {font=\small},
  hlbl/.style  = {font=\small\bfseries},
  dashline/.style = {dashed, thick, gray!50},
]

% 生命线
\draw[thick] (0, 0) -- (0, -8.5);
\draw[thick] (6, 0) -- (6, -8.5);
\draw[thick] (12,0) -- (12,-8.5);

% 角色标题
\node[hlbl, above=0.1cm] at (0,  0) {客户端（浏览器）};
\node[hlbl, above=0.1cm] at (6,  0) {服务器};
\node[hlbl, above=0.1cm] at (12, 0) {订单房间};

% Phase 1: HTTP 升级握手
\draw[dashline] (-1.5,-1.0) -- (13.5,-1.0);
\node[font=\small\itshape, right] at (-1.5,-0.8) {阶段一：HTTP 升级握手};

\draw[arr] (0,-1.4) -- node[above, lbl]{HTTP GET, Upgrade: websocket} (6,-1.4);
\draw[arr] (6,-2.0) -- node[above, lbl]{101 Switching Protocols} (0,-2.0);

% Phase 2: 身份认证
\draw[dashline] (-1.5,-2.7) -- (13.5,-2.7);
\node[font=\small\itshape, right] at (-1.5,-2.5) {阶段二：Token 认证};

\draw[arr] (0,-3.1) -- node[above, lbl]{\{event: join, orderID, token\}} (6,-3.1);
\draw[arr] (6,-3.7) -- node[above, lbl]{验证 Token → 通过} (0,-3.7);
\draw[arr] (6,-4.2) -- node[above, lbl]{加入 Room(orderID)} (12,-4.2);

% Phase 3: 消息交互
\draw[dashline] (-1.5,-4.9) -- (13.5,-4.9);
\node[font=\small\itshape, right] at (-1.5,-4.7) {阶段三：全双工消息交互};

\draw[arr] (0,-5.3) -- node[above, lbl]{发送消息帧（Data Frame）} (6,-5.3);
\draw[arr] (6,-5.8) -- node[above, lbl]{加密存储 + 广播到 Room} (12,-5.8);
\draw[arr] (12,-6.3)-- node[above, lbl]{推送给房间内所有成员} (0,-6.3);

% Phase 4: 连接关闭
\draw[dashline] (-1.5,-7.0) -- (13.5,-7.0);
\node[font=\small\itshape, right] at (-1.5,-6.8) {阶段四：连接关闭};

\draw[arr] (0,-7.4) -- node[above, lbl]{Close Frame} (6,-7.4);
\draw[arr] (6,-8.0) -- node[above, lbl]{Close ACK + 触发 Merkle 快照} (0,-8.0);

\end{tikzpicture}
\caption{WebSocket 握手及订单通信时序示意图}
\label{fig:websocket_seq}
\end{figure}


需要指出的是，WebSocket 协议本身主要解决实时通信的连接管理问题，并不直接提供消息加密、身份认证或日志完整性保护等安全功能。因此，在实际应用中，WebSocket 应结合 TLS 形成 WSS（WebSocket Secure）安全连接，以防止通信内容在传输过程中被窃听或篡改。本作品在系统设计中将 WebSocket 作为实时消息通道，将 Token 认证、消息加密、哈希审计和权限控制作为安全增强机制，从而避免将普通聊天功能误认为完整安全机制。
